\documentclass[a4paper,9pt]{extarticle}

%=========================================================
% PACKAGES
%=========================================================

\usepackage[
a4paper,
top=1.8cm,
bottom=1.5cm,
footskip=7mm,
left=1cm,
right=1cm,
headheight=14pt,
headsep=5mm
]{geometry}

\usepackage[T1]{fontenc}
\usepackage[utf8]{inputenc}
\usepackage[french]{babel}

\usepackage[sfdefault]{FiraSans}
\renewcommand{\familydefault}{\sfdefault}

\usepackage{multicol}
\setlength{\columnsep}{0.7cm}

\usepackage{amsmath,amssymb}
\usepackage{titlesec}
\usepackage{enumitem}

\usepackage{tcolorbox}
\tcbuselibrary{skins,breakable}

\usepackage{xcolor}
\usepackage{graphicx}
\usepackage{float}

\usepackage{fancyhdr}
\usepackage[hidelinks]{hyperref}


%=========================================================
% PAGE SETTINGS
%=========================================================

\setlength{\parindent}{0pt}
\setlength{\parskip}{1.5pt}

\setlength{\abovedisplayskip}{4pt}
\setlength{\belowdisplayskip}{4pt}

\setlist[itemize]{
leftmargin=*,
itemsep=1pt,
topsep=1pt
}
\setlist[enumerate]{
leftmargin=*,
itemsep=1pt,
topsep=1pt
}


%=========================================================
% COLORS
%=========================================================

\definecolor{Blue}{HTML}{E3F0FF}
\definecolor{Green}{HTML}{E3F7EA}
\definecolor{Pink}{HTML}{FFE5ED}
\definecolor{Orange}{HTML}{FFF0DC}
\definecolor{Purple}{HTML}{EEE7FF}
\definecolor{Yellow}{HTML}{FFFADB}
\definecolor{Goat}{HTML}{F7EEDF}

\definecolor{BlueTitle}{HTML}{285B8F}
\definecolor{GreenTitle}{HTML}{287052}
\definecolor{PinkTitle}{HTML}{A84063}
\definecolor{OrangeTitle}{HTML}{A96500}
\definecolor{PurpleTitle}{HTML}{6245A5}
\definecolor{GoatTitle}{HTML}{8A5A1F}
\definecolor{RecapTitle}{HTML}{5B3A8E}


%=========================================================
% HEADER
%=========================================================

\pagestyle{fancy}
\fancyhf{}

\fancyhead[L]{\small\textbf{Comportement mécanique des matériaux}}
\fancyhead[C]{\small Compression, dureté, fluage, rupture, fatigue}
\fancyhead[R]{\small Week 02 --- 14-20/09/2026}

\fancyfoot[C]{\small\href{https://github.com/AlphaXT38G/notes_ba5}{github.com/AlphaXT38G/notes\_ba5}}
\fancyfoot[R]{\small \thepage}

\renewcommand{\headrulewidth}{0.4pt}


%=========================================================
% TITLES
%=========================================================

\titleformat{\section}
{\large\bfseries\color{BlueTitle}}
{}{0pt}{}

\titlespacing{\section}
{0pt}{6pt}{3pt}


%=========================================================
% BOXES
%=========================================================
% Boxes are kept whole (unbreakable): inside multicols, a box split by tcolorbox
% leaves an empty column behind it.

\newtcolorbox{definition}{
enhanced, unbreakable, colback=Blue, colframe=BlueTitle,
boxrule=0.7pt, arc=0mm, title=\textbf{Définition}, fonttitle=\bfseries
}

\newtcolorbox{important}{
enhanced, unbreakable, colback=Pink, colframe=PinkTitle,
boxrule=0.7pt, arc=0mm, title=\textbf{Idée clé}, fonttitle=\bfseries
}

\newtcolorbox{formula}{
enhanced, unbreakable, colback=Green, colframe=GreenTitle,
boxrule=0.7pt, arc=0mm, title=\textbf{Formules clés}, fonttitle=\bfseries
}

\newtcolorbox{warning}{
enhanced, unbreakable, colback=Orange, colframe=OrangeTitle,
boxrule=0.7pt, arc=0mm, title=\textbf{Attention}, fonttitle=\bfseries
}

\newtcolorbox{example}{
enhanced, unbreakable, colback=Purple, colframe=PurpleTitle,
boxrule=0.7pt, arc=0mm, title=\textbf{Exemple}, fonttitle=\bfseries
}

% GOAT - brown: the idea told as simply as possible
\newtcolorbox{goat}{
enhanced, unbreakable, colback=Goat, colframe=GoatTitle,
boxrule=0.7pt, arc=0mm, title=\textbf{Expliqué à une chèvre}, fonttitle=\bfseries
}

% RECAP - violet: what to know for the exam, at the end of the week
\newtcolorbox{recap}{
enhanced, breakable, colback=Purple, colframe=RecapTitle,
boxrule=1pt, arc=0mm, title=\textbf{Récap pour l'examen}, fonttitle=\bfseries
}

% Screenshot of a slide of the course: \slide{cNN-PP}
\newcommand{\slide}[2][0.95]{%
\begin{center}
\includegraphics[width=#1\linewidth]{../images/course/#2}
\end{center}}


%=========================================================
% DOCUMENT
%=========================================================

\begin{document}

\thispagestyle{fancy}

\begin{multicols}{2}

%=========================================================
\section{Compression, dureté, flexion}

\begin{important}
\textbf{Compression} : deux défauts opposés. \textbf{Flambage} $\Rightarrow$ barreaux trapus ;
\textbf{frottement} (mise en tonneau) $\Rightarrow$ barreaux longs. Compromis
$h/d \approx 3$--$4$. Courbe conventionnelle plus haute qu'en traction (section qui grandit, pas
de striction). Utile si pas de barreau de traction possible, ou béton, composites.
\end{important}

\slide[0.75]{c01-47}

\begin{formula}
\[ H = \frac{F}{A} \qquad\text{et}\qquad H \approx 3\sigma_y \]
(sans écrouissage, déformation 2D). Le métal sous l'indenteur est confiné : il faut $\approx 3$
fois $\sigma_y$ pour le faire couler.
\end{formula}

Essais : Brinell (bille), Vickers et Knoop (pyramides), Rockwell (profondeur, table),
instrumenté (nanoindentation, courbe force--enfoncement). Unités conventionnelles (HV\ldots).

\begin{formula}
Flexion (Euler--Bernoulli) : $\varepsilon_{xx} = y/R$ ($y$ : distance à l'axe neutre,
$R$ : courbure). Profil $\sigma(y)$ = courbe $\sigma(\varepsilon)$ à l'échelle $1/R$ près.
\end{formula}

\begin{warning}
Remonter à $\sigma(\varepsilon)$ à partir de ($P$, $\delta$) en flexion n'est pas simple.
Flexion 3 points (moment max au centre) / 4 points (moment constant au milieu).
\end{warning}

\begin{goat}
Écraser : trop long ça se plie, trop court ça frotte et gonfle. Dureté : on enfonce une bille, la
pression dans le trou vaut trois fois la résistance en traction. Plier : une face tirée, l'autre
poussée.
\end{goat}

%=========================================================
\section{Quand le temps compte}

\begin{important}
À $T > \frac12 T_f$ (K) pour métaux et céramiques (dès l'ambiante pour les polymères), la
déformation dépend du temps : \textbf{vitesse}, \textbf{relaxation} (déformation fixée,
$\sigma$ baisse), \textbf{fluage} (charge fixée, $\varepsilon$ monte).
\end{important}

\begin{formula}
\[ \sigma = K\dot\varepsilon^{\,m}, \quad m = \frac{\partial\ln\sigma}{\partial\ln\dot\varepsilon}
   \approx \frac{\ln(\sigma_2/\sigma_1)}{\ln(\dot\varepsilon_2/\dot\varepsilon_1)} \]
(essai de saut de vitesse).
\[ \dot\varepsilon = A\,\sigma^N e^{-Q/RT} \]
fluage stationnaire ; $N$ entre 1 et 10 (5 : dislocations), $Q \sim 10^5$\,J/mol ;
$\mathbf{m = 1/N}$.
\end{formula}

\slide[0.85]{c01-59}

Fluage : primaire (ralentit), \textbf{secondaire} (constant, le plus long), tertiaire (accélère
jusqu'à rupture). $N$ = pente de $\log\dot\varepsilon$ vs $\log\sigma$ ; $-Q/R$ = pente de
$\ln\dot\varepsilon$ vs $1/T$.

\begin{important}
$m$ élevé freine la striction (superplasticité). $m = 1$ : \textbf{fluide newtonien}
$\sigma = K\dot\varepsilon$, \textbf{pas de striction} (verre chaud, caramel).
\end{important}

\begin{goat}
À chaud, la matière coule comme du miel épais : sous un poids elle s'allonge peu à peu, et si
elle durcit quand on la tire plus vite, aucun endroit ne s'amincit en avance.
\end{goat}

%=========================================================
\section{Mécanique de la rupture}

\begin{definition}
Mécanique des \textbf{corps fissurés}. Deux paramètres :
\begin{itemize}
\item \textbf{énergie} : $G = \pi F(a)\sigma^2a/E$ (force motrice), rupture si $G = G_c$
      (énergie pour créer 1\,m$^2$ de fissure, J/m$^2$) ;
\item \textbf{contraintes} : en pointe, $\sigma \propto K/\sqrt{2\pi r}$ ; rupture si $K = K_c$
      (ténacité, MPa$\sqrt{\text m}$).
\end{itemize}
\end{definition}

\begin{formula}
\[ K = Y\sigma\sqrt{\pi a}, \qquad \sigma_r = \frac{K_c}{Y\sqrt{\pi a}}, \qquad G_c = A\frac{K_c^2}{E} \]
$Y = 1$ : fissure interne de \textbf{demi}-longueur $a$, plaque infinie ; $Y = 1{,}12$ : fissure
de bord de longueur $a$.
\end{formula}

\slide[0.9]{c01-69}

\begin{important}
$K$ croît avec $a$ : une fois partie, la fissure ne s'arrête pas. $K = K_c$ = rupture
\textbf{catastrophique}. On travaille avec $K$ car il est proportionnel à $\sigma$.
\end{important}

\begin{example}
\textit{(pas dans les diapos)} Verre, $K_c = 0{,}7$, rayure de bord $a = 10\,\mu$m :
$\sigma_r = 0{,}7/(1{,}12\sqrt{\pi\cdot10^{-5}}) \approx 110$\,MPa.
\end{example}

\begin{goat}
Toute la force se concentre au bout d'une déchirure. $K$ dit à quel point ; quand il atteint ce
que le matériau supporte, la déchirure file et ne s'arrête plus.
\end{goat}

%=========================================================
\section{Essais de rupture et fatigue}

\begin{itemize}
\item \textbf{Ténacité} : éprouvette entaillée + préfissure, charge de départ $\to K_c$.
\item \textbf{Charpy} (Izod) : pendule, énergie absorbée ; montre la \textbf{transition
      ductile--fragile} des aciers.
\item \textbf{Formabilité} des tôles : pliage, emboutissage (Erichsen, Swift\ldots), cornes.
\end{itemize}

\slide[0.85]{c01-80}

\begin{important}
\textbf{Fatigue}, charge cyclique, rupture bien sous $R_m$.
(i) Éprouvette lisse : \textbf{Wöhler} (S--N), $N_f$ à amplitude donnée ;
\textbf{statistique}, sensible à la fréquence (échauffement des polymères), l'environnement, la
surface.
(ii) Éprouvette fissurée : $da/dN$ en fonction de $\Delta K$.
\end{important}

\begin{formula}
\[ \Delta K = K_{\max} - K_{\min} = Y\Delta\sigma\sqrt{\pi a}, \qquad \frac{da}{dN} = A(\Delta K)^m \]
($K_{\min} = 0$ si $\sigma_{\min} < 0$). Seuil $\Delta K_{th}$ ; rupture quand $K_{\max} = K_c$.
\end{formula}

\slide[0.75]{c01-84}

\begin{goat}
Plier un trombone vingt fois : à chaque aller-retour la petite fissure avance d'un pas, plus
grand si elle est déjà longue. Quand elle est assez longue, tout casse.
\end{goat}

%=========================================================
\section{Récap de la semaine}

\begin{recap}
\begin{itemize}
\item Compression : flambage vs tonneau, $h/d \approx 3$--$4$, pas de striction.
\item $H = F/A \approx 3\sigma_y$. Flexion : $\varepsilon_{xx} = y/R$.
\item $\sigma = K\dot\varepsilon^m$ ; relaxation ; fluage 3 stades ;
      $\dot\varepsilon = A\sigma^Ne^{-Q/RT}$ ; $m = 1/N$ ; newtonien $m = 1$ sans striction.
\item $K = Y\sigma\sqrt{\pi a}$, rupture à $K_c$ ; $Y = 1$ (interne, $a$ demi-longueur) ou
      $1{,}12$ (bord) ; $G_c = AK_c^2/E$ ; rupture catastrophique.
\item Charpy, ténacité, formabilité ; Wöhler ; $da/dN = A(\Delta K)^m$.
\end{itemize}
\end{recap}

\end{multicols}

\end{document}
